\chapter{طبقة سجل العمليات (Logging Layer)}
\label{CH:LOG}

واحدة من أكبر المشاكل في تصميم أنظمة الملفات هي كيفية التعامل مع الانقطاع المفاجئ للطاقة أو انهيار النظام أثناء كتابة البيانات على القرص الصلب.
لأن العديد من عمليات نظام الملفات تتطلب تعديل عدة كتل مختلفة على القرص (مثل تخصيص كتلة جديدة، وتعديل خريطة البتات، وتحديث العقدة الفهرسية، وإضافة مدخل في الدليل).

إذا انقطعت الكهرباء في منتصف الطريق، أصبحت حالة القرص الصلب إما متهالكة أو مفقودة التناسق (Inconsistent State).

يحل \texttt{xv6} هذه المشكلة عبر تقديم \textbf{طبقة سجل العمليات} (Logging Layer / Journaling).

يقدم هذا الفصل ميكانيكية عمل سجل العمليات، والمعاملات الذرية (Atomic Transactions)، وكيفية التعافي عند تشغيل النظام.

\section{غاية السجل وتصميمه}
\label{sec:log_design}

يقتضي مفهوم سجل العمليات عدم كتابة التعديلات مباشرةً على كتل القرص الأصلية.
بدلاً من ذلك، تتبع النواة الخطوات التالية:
\begin{enumerate}
  \item تُسجل كافة الكتل المعدلة في منطقة معزولة خاصة على القرص تُدعى \textbf{منطقة السجل} (Log Area).
  \item عند اكتمال جمع جميع كتل المعاملة، تكتب النواة كتلة رأس السجل (Header Block) على القرص. تُعرف هذه الخطوة بـ \textbf{الالتزام} (Commit).
  \item بعد تأكيد الالتزام، تقوم النواة بنسخ التعديلات من منطقة السجل إلى مواقعها الفيزيائية الأصلية على القرص. تُسمى هذه الخطوة بـ \textbf{التفريغ} (Install).
  \item بعد تفريغ كل الكتل، تصفح النواة السجل لتفريغه للعمليات الجديدة.
\end{enumerate">

إذا انهارت النواة أو انقطعت الكهرباء قبل الالتزام (قبل كتابة كتلة الرأس)، يتم تجاهل السجل بالكامل عند الإقلاع وتظل البيانات القديمة كما هي.
وإذا انهارت النواة بعد الالتزام وقبل إكتمال التفريغ، يقرأ النظام عند الإقلاع كتلة الرأس ويتولى إعادة تفريغ الكتل من السجل إلى مواقعها الأصلية بنجاح، مما يضمن الذرية التامة للمعاملات.

\section{الكود البرمجي: تنفيذ السجل}
\label{sec:code_logging}

تُدار عمليات السجل في \texttt{kernel/log.c}.
تغلف النواة أية عملية تعديل على نظام الملفات بين الاستدعائين \texttt{begin\_op()} و \texttt{end\_op()}:
\begin{lstlisting}
begin_op();
// تعديل العقد الفهرسية والكتل المخصصة عبر log_write()
ilock(ip);
ip->size = new_size;
iupdate(ip);
iunlock(ip);
end_op();
\end{lstlisting}

تبدأ الدالة \texttt{begin\_op()} بالانتظار حتى لا يكون السجل خاضعًا للتثبيت أو ممتلئًا بمعاملات جارية:
\begin{itemize}
  \item تزيد عداد المعاملات المشتركة الجارية \texttt{log.outstanding}.
  \item عندما تدعو العمليات \texttt{log\_write(b)}، لا تُكتب الكتلة على القرص فورًا؛ بل تُضاف إلى مصفوفة السجل في الذاكرة وتسجل في الهيكل.
\end{itemize}

عند دعوة \texttt{end\_op()}:
\begin{enumerate}
  \item ينقص عداد المعاملات الجارية.
  \item إذا وصل العداد إلى الصفر وكانت هناك كتل معلقة، تدعو النواة الدالة \texttt{commit()}.
  \item تقوم \texttt{commit()} بدعوة \texttt{write\_log()} لكتابة الكتل المعدلة في منطقة السجل على القرص.
  \item تدعو \texttt{write\_head()} لكتابة كتلة رأس السجل على القرص (نقطة الذرية والالتزام).
  \item تدعو \texttt{install\_trans()} بنسخ الكتل من منطقة السجل إلى أماكنها المخصصة في القرص.
  \item تصفر كتلة الرأس وتدعو \texttt{write\_head()} لتفرغ السجل المعاملاتي.
\end{enumerate}

\section{العالم الحقيقي}
\label{sec:real_world_log}

تعتمد أنظمة الملفات الإنتاجية مثل ext4 و NTFS و XFS تقنيات سجل الميتاداتا (Metadata-only Logging) حيث تُسجل البيانات الوصفية فقط دون كتل البيانات العادية لتوفير الأداء والسرعة.

\section{تمارين}
\label{sec:exercises_log}

\begin{enumerate}
  \item اشرح ماذا يحدث لو انقطعت الكهرباء أثناء تنفيذ الدالة \texttt{install\_trans()} وكيف يتم التعافي عند الإقلاع القادم.
  \item قم بزيادة حجم السجل واختبر تأثيره على سرعة المعاملات المتوازية في \texttt{xv6}.
\end{enumerate}
